\documentclass[10pt]{article}
\usepackage[letterpaper,margin=0.8in]{geometry}
\usepackage{amsmath,amssymb,bm}
\usepackage[hidelinks]{hyperref}
\setlength{\parindent}{0pt}
\setlength{\parskip}{4pt}
\title{Decoupling Return Seeking and Risk Control\\in Portfolio Reinforcement Learning}
\author{Name (Student ID)\\Department, Seoul National University}
\date{}
\begin{document}
\maketitle
\vspace{-1.5em}

\section{Motivation}
Portfolio RL often combines return and risk in one weighted reward, obscuring whether improvements arise from forecasting, sequential control, or simply reducing market exposure. We ask: \emph{Can a separately trained risk controller improve the out-of-sample return--risk frontier of a frozen return-seeking policy, beyond conventional risk overlays?} Inspired by performance--safety decoupling in ResSafe~\cite{ressafe}, we study a financial adaptation with explicit portfolio accounting and transaction costs. This is an empirical hypothesis, not a claim of guaranteed safety or a first use of risk-aware portfolio RL.

\section{Problem Definition}
Consider a small price-taking investor allocating among $K=3$--$5$ liquid Korean ETFs and cash, without leverage or short sales. Let $x_t$ summarize observable market history and let $z_t$ contain holdings, cash, wealth, and running peak wealth. Actions specify target portfolio weights $w_t\in\Delta^{K+1}$. Market dynamics are exogenous, but account dynamics are controlled:
\begin{equation}
 P(x_{t+1}\mid x_{\leq t},w_t)=P(x_{t+1}\mid x_{\leq t}),\qquad
 z_{t+1}=F(z_t,w_t,x_{t+1}).
\end{equation}
Here $F$ includes execution prices, drifted holdings, transaction costs, and cash feasibility. A finite history is an approximate information state, not an assumed exact Markov state. Define net log growth $g_{t+1}=\log(V_{t+1}/V_t)$ and drawdown $d_t=1-V_t/\max_{u\leq t}V_u$. Close-of-day information generates orders for the next trading day's open; neither observations nor fills may use unavailable future data.

\section{Literature Review}
ResSafe freezes a task policy and trains a residual action correction~\cite{ressafe}; financial risk requires a different objective and evaluation. Lexicographic RL provides a framework for prioritizing objectives~\cite{lex}, and risk-constrained RL treats tail risk through explicit constraints~\cite{chow}. Neural-progressive hedging already combines RL with model-guided portfolio constraint enforcement~\cite{np}. Thus, modular risk control alone is insufficient novelty. Our focus is whether exploiting exogenous market dynamics enables a useful, reusable controller under trading frictions and changing market regimes.

\section{Main ideas and novelty}
\textbf{Two-stage training.} First train a nominal PPO allocation policy to maximize $J(\pi)=\mathbb{E}[T^{-1}\sum_t g_{t+1}]$. Freeze its parameters, then train a residual controller receiving the current information state and nominal allocation $w_t^0$. Use
\begin{equation}
 w_t=\operatorname{Proj}_{\Delta^{K+1}}\!\left[w_t^0+\delta_\phi(x_t,z_t,w_t^0)\right].
\end{equation}
Actual holdings must evolve under the corrected actions; the nominal policy is queried on the resulting state. For estimated conditional covariance $\widehat\Sigma_t$ (including cash), the primary second-stage problem is
\begin{equation}
 \min_\phi\;\mathbb{E}\!\left[\frac1T\sum_t
 \left(w_t^\top\widehat\Sigma_t w_t+\eta\|w_t-w_t^0\|_2^2\right)\right]
 \quad\text{s.t.}\quad J(\pi_\phi)\geq J(\pi_0)-\epsilon.
\end{equation}
The tolerance $\epsilon$ is fixed using training/validation data and measured in mean daily net log-return units. An adaptive multiplier approximately enforces this empirical return constraint. It offers no out-of-sample guarantee. Separate return/risk critics estimate the signals; they do not by themselves solve the trade-off. Downside squared loss and drawdown-increment penalties are controlled reward ablations, rather than an unrestricted reward search. Variance reduction alone can collapse to cash, so return retention and exposure-matched baselines are essential.

\textbf{Exogenous scenarios and model-based ablation.} Compare model-free training on historical episodes against model-based residual training on short scenarios from a fitted VAR model with jointly resampled residuals; a small probabilistic GRU is optional. Only the market process is learned; account transitions remain explicit. The same market paths evaluate different allocations with common random numbers under the price-taking simulator. This reduces comparison noise, but does not create new independent market evidence. AR/GRU prediction alone is not model-based RL: generated transitions must enter policy optimization. Test whether benefits survive forecast misspecification, rather than assuming a larger predictor helps.

\section{Expected contributions}
The intended contributions are a structured formulation separating uncontrolled market dynamics from controlled account dynamics; a return-preserving residual-risk algorithm; and a controlled empirical study of when sequential learning helps beyond one-step allocation and standard risk overlays. As a limited theoretical result, establish that with action-independent market observations, no trading costs, no account-dependent constraints, and additive per-period utility depending only on current market information and allocation, optimization reduces to conditional one-step utility maximization. Costs and path-dependent objectives motivate the sequential setting studied here.

\section{Detailed plan}
\textbf{Week 1: feasibility and simulator.} Confirm IntelliQuant history entitlements and permission/support for training-data export. Public documentation describes JavaScript strategies, historical market/fundamental access, execution settings, and paid-tier result downloads~\cite{iq}; a Python training API and raw-data bulk export remain unverified. Train locally only on legitimately accessible data. Use IntelliQuant for final backtest replication via a compact JavaScript policy, subject to runtime limits. If data extraction is unavailable, use a separately licensed dataset for training and disclose the source difference. Verify next-open execution, distributions/adjustments, costs, missing data, and ETF listing dates; do not choose today's winners retrospectively.

\textbf{Weeks 2--3: core experiments.} Implement the nominal policy, joint return--risk PPO, and residual PPO. Compare with equal weight, nominal plus volatility targeting, nominal plus a covariance-based minimum-deviation optimization filter, and a forecast-plus-one-step allocator. If fractional-share accounting is differentiable, include direct policy optimization through historical replay as a sequential baseline. Match information, allocation limits, costs, and tuning budgets. Fit scalers, covariance estimators, and any predictors using available past data only. Freeze all learned parameters within each test block.

\textbf{Weeks 4--5: validation and extensions.} Run chronological walk-forward evaluation with distinct training, validation, and untouched test blocks; purge overlapping target horizons. Add the VAR scenario experiment, then the GRU only if time permits. Report net growth, volatility, Sharpe, maximum drawdown, turnover, realized tail losses, return-constraint violations, and average cash exposure. Use multiple training seeds and paired moving-block bootstrap intervals; seeds are not independent market samples. Compare at matched volatility/exposure, test higher costs and adverse regimes, and report both training stages' compute. Removing costs and path dependence tests whether RL's advantage disappears. A null result against conventional controls is an informative outcome.

\clearpage
\begin{thebibliography}{9}
\bibitem{ressafe} G. Qu et al. \emph{ResSafe: Learning Safety Filtering with Residual Reinforcement Learning for Humanoids}. arXiv:2609.15988, 2026. \url{https://arxiv.org/abs/2609.15988}.
\bibitem{lex} J. Skalse, L. Hammond, C. Griffin, and A. Abate. \emph{Lexicographic Multi-Objective Reinforcement Learning}. IJCAI, pp. 3430--3436, 2022. \url{https://www.ijcai.org/proceedings/2022/476}.
\bibitem{chow} Y. Chow, M. Ghavamzadeh, L. Janson, and M. Pavone. \emph{Risk-Constrained Reinforcement Learning with Percentile Risk Criteria}. JMLR, 18(167):1--51, 2018. \url{https://www.jmlr.org/papers/v18/15-636.html}.
\bibitem{np} S. Ghosh, L. Wynter, S. H. Lim, and D. T. Nguyen. \emph{Neural-progressive hedging: Enforcing constraints in reinforcement learning with stochastic programming}. UAI, PMLR 180:707--717, 2022. \url{https://proceedings.mlr.press/v180/ghosh22a.html}.
\bibitem{iq} IntelliQuant. \emph{Strategy/API documentation and membership}. Accessed October 28, 2026. \url{https://www.intelliquant.ai/help/manual/1}; \url{https://www.intelliquant.ai/help/ref/15}; \url{https://www.intelliquant.ai/membership}.
\end{thebibliography}
\end{document}
